A grouping can make measurement comparable or description informative.
The distinctions among \mention{damn} uses have descriptive value even
before a new prediction is tested. \term{Projectibility} concerns the
warrant for expectations beyond the evidence used to identify a case
\citep{Goodman1955}. Neither descriptive usefulness nor a successful
projection identifies its own causal explanation. These achievements
need separate assessment.

Within linguistics, \textcite[663, 666, 680]{haspelmath2010} distinguishes
language-particular descriptive categories from concepts designed for
cross-language comparison. The division of labour here is different:
several inquiries within English can organize the same expressions
differently. The assessment developed here identifies where evidence
for one claim leaves another unsupported. A semantic contrast can be
preserved without establishing category-wide projections; phonological
resemblance can leave grammatical membership unresolved. Specifying
the bearer, grouping, and intended inference makes those transitions
available for adjudication.

For the selected \mention{Bang!}, the relevant contribution is a portrayal
of an impact sound. For \mention{Damn!}, it's the conveyed stance toward
a situation. The distinction changes the interpretation of a possible
modification. In a constructed narrative, \mention{bang, bang, bang} can
multiply portrayed impacts; repeated \mention{damn} can strengthen the
conveyed stance. Repetition can thus transform the contribution in
different ways. These examples articulate the contrast to be explained;
they aren't results from a new elicitation study.

Descriptive success can be assessed without requiring a prediction about
a differently named property. The expressive grouping preserves a semantic relation across
standalone and integrated constructions; the expressive--depictive contrast
records different contributions concealed by the broad cover term.
Those are warranted uses of descriptive classes. Calling them descriptions
doesn't diminish their scientific value. It specifies the achievement
so that categoryhood, in the stronger sense used here, remains an
answerable further question.

\begin{table}[htbp]
\centering
\small
\caption{Assessing the expressive grouping without conflating its achievements}
\label{tab:inquiry-record}
\begin{tabular}{@{}>{\raggedright\arraybackslash}p{.17\textwidth}
                  >{\raggedright\arraybackslash}p{.36\textwidth}
                  >{\raggedright\arraybackslash}p{.39\textwidth}@{}}
\toprule
Question & Proposed answer & Evidential consequence \\
\midrule
Bearer & Conventional uses, including standalone and prenominal \mention{damn} &
Identify each use without requiring the target property under investigation. \\
\addlinespace
Grouping & Deployment conventionally conveys evaluative or affective stance &
Distinguish this from describing an emotion, depicting a scene, or managing interaction. \\
\addlinespace
Contribution & Preserve expressive commonality across grammatical constructions &
Distinguish the observed contrast from its extension to new realizations, environments, or sign-types. \\
\addlinespace
Alternative & An undifferentiated interjection class, explicit subtypes, or item-specific accounts &
Compare which semantic relations each preserves and what further generalization each warrants. \\
\bottomrule
\end{tabular}
\end{table}

An admission rule fixes which cases enter a class. A category claim
makes a further commitment: membership supports informative expectations
beyond the evidence used to assign it. Assembling a grouping and
warranting those expectations can succeed or fail separately. Here,
category standing names the latter achievement; it isn't a verdict
on whether a descriptive class is scientifically useful.

The reach of an expectation needs specifying. Evidence about a familiar
sign can warrant expectations about new realizations or environments.
Extending an expressive contrast to another construction can already
involve induction. This establishes reach for that sign without by
itself warranting the same expectation for another sign-type. A
category-level projection can make that further commitment about an
independently identified member.

The target needn't be a differently named property. A familiar property
in a new environment can provide informative reach if it wasn't part
of the admission evidence. Assigning expressive uses by their stance
in one construction and investigating how that contribution behaves
elsewhere differs from defining the class to contain only uses already
known to exhibit the target behavior.

An additional shared property isn't sufficient merely because it's
additional. If all members and nearly every sign outside a proposed
grouping exhibit the target repetition behavior, membership supplies
little additional guidance about it. An imperfect probabilistic association can be
informative when it changes a warranted expectation at the stated
scope. The alternatives in the preceding section help assess that
contribution.

This comparison concerns what membership licenses an investigator to
expect. Whether a category representation improves a particular model
is a separate question, taken up below. Projective warrant doesn't
require that representation to outperform every account of the same
evidence. Both claims remain limited to the profile and scope tested:
support for a subtype doesn't automatically extend to its superordinate
class.

I provisionally accept the grammatical and action-specific categories
used here for the relations documented in the cited descriptions.
This inherits their evidential warrant at that scope; it doesn't
establish every extension considered here. The present comparisons
defend expressive and depictive distinctions and reject an unsupported
transfer from phonology to grammar. They leave the broader projective
standing of the proposed semantic, phonological, and ancestry groupings
open.

Consider a hypothetical adverse case for standalone stance \mention{damn}.
The initial identification uses its conventional reaction meaning and
recurrence across independently sampled situations; the grammatical
prediction concerns resistance to integration. A new occurrence retains
that reaction meaning but functions as an integrated dependent of a
larger constituent. The first response is to check that structural
analysis. If it stands, the occurrence counts against the proposed
distribution under the original identification rule. Its inconvenience
isn't evidence for a second sign.

Evaluate the original projection under the identification rule available
when it was made. Later evidence may warrant distinguishing a recurrent
conventional use by its contribution or distribution. That revision
can improve subsequent predictions, but it doesn't remove earlier
counterexamples from the original assessment. Observations that motivate
a new distinction don't also supply independent confirmation of its
predictive success.

A quoted token or a productive noun denoting an utterance remains
another possible analysis, but it needs evidence for quotation or
nominal reference. A correction to a mistaken structural analysis should
be recorded as such. It differs from changing the identification rule
so that a confirmed counterexample no longer falls within its scope.

\textcite[295--299]{hill1965} offered viewpoints for comparing causal
interpretations with alternatives, including temporal order and consistency
with independently supported knowledge. They concern specified causes
and outcomes; they neither establish category standing nor form a
mechanical checklist. Section~\ref{sec:history} applies the temporal
question to \mention{gee}. Recurrence alone identifies neither a causal
ordering nor a process that maintains a profile.

The grammatical claim concerns membership in a lexical category. The
evidence offered concerns sounds or sound patterns unusual elsewhere in
the language. That evidence supports a comparison of phonological
possibilities across independently identified signs. By itself, it doesn't
settle whether a particular use has nominal, verbal, or interjectional
grammar. Nor does sound depiction itself entail phonological anomaly: resemblance
and distribution within the sound system are different relations.

Phonological evidence can support a grammatical expectation where an
independently established association warrants it. The shared unusualness
considered here doesn't itself establish that association. The limitation
concerns the evidence connecting the properties, not the fact that they
belong to different domains.

For English \mention{Bang!}, the competing grammatical analysis is a
standalone nominal fragment, supported as a possibility by the independently
available noun in \mention{a loud bang}. A depictive interpretation leaves
that analysis open. Even if an occurrence had unusual phonation, the
additional observation would concern its realization. Choosing between
nominal and interjectional analyses needs distributional and constructional
evidence about the selected conventional use, or an independently warranted
association connecting the observed phonology to that grammar.

Hill's considerations clarify a tempting stronger claim. If newly acquired
prosodic isolation were proposed as a cause of the semantic change, the
already peripheral invocation would fail to supply the required change
in exposure. The cited history doesn't independently measure a prosodic
innovation preceding the semantic transition. Plausibility of a
prosody--syntax connection can't supply that missing temporal comparison.
An already present peripheral construction could still be an enabling
condition for later change.
Gehweiler's proposed roles for invited inference and taboo-motivated
clipping remain distinct historical claims with their own evidence.

The \href{inquiry-relative-evidence-supplement.pdf}{evidence supplement}
preserves the complete records: GloWbE measurement and model comparisons
(Section~A), the derivative audit (Section~B), the \mention{fie} audit
(Section~C), and inventory turnover (Section~D). Its retained model
comparison predicts an omitted country--item cell within the observed
design; it doesn't assess new countries, new items, or the proposed
linguistic groupings. These records locate
the failure at the relevant claim. A useful retrieval class survives
as a tool, but its output can't establish a linguistic category
without evidence connecting that class to independently identified uses.
